\documentclass[11pt,a4paper]{article}

% ---------- encoding & fonts ----------
\usepackage[utf8]{inputenc}
\usepackage[T1]{fontenc}
\usepackage{lmodern}
\usepackage[spanish,es-tabla]{babel}

% ---------- math ----------
\usepackage{amsmath,amssymb,amsthm}

% ---------- graphics & diagrams ----------
\usepackage{tikz}
\usetikzlibrary{positioning,arrows.meta,shapes,fit,backgrounds,calc,babel}
\usepackage{pgfplots}
\pgfplotsset{compat=1.18}

% ---------- tables ----------
\usepackage{booktabs,multirow}

% ---------- hyperlinks ----------
\usepackage[colorlinks=true,linkcolor=blue,citecolor=blue,urlcolor=blue]{hyperref}

% ---------- geometry ----------
\usepackage[top=2.5cm,bottom=2.5cm,left=2.5cm,right=2.5cm]{geometry}

% ---------- theorems ----------
\theoremstyle{definition}
\newtheorem{definicion}{Definición}[section]
\newtheorem{teorema}{Teorema}[section]
\newtheorem{observacion}{Observación}[section]

% ---------- custom commands ----------
\newcommand{\Risk}{\mathcal{R}}
\newcommand{\Prob}{\mathbb{P}}
\newcommand{\E}{\mathbb{E}}
\newcommand{\X}{\mathcal{X}}
\newcommand{\A}{\mathcal{A}}
\newcommand{\R}{\mathbb{R}}
\newcommand{\Beta}{\operatorname{Beta}}

% ---------- title ----------
\title{\textbf{Agente Estadístico de Wald para Decisiones Crediticias}\\{\large Implementación en LangGraph de un meta-selector Bayes--Minimax}}
\author{Lab F.7 — Capítulo 7}
\date{\today}

\begin{document}
\maketitle

\tableofcontents
\newpage

% ===================================================================
\section{Introducción}
% ===================================================================

Este documento describe un agente de decisión estadística que implementa la
\textbf{teoría de Wald} del Capítulo~7 para la aprobación de solicitudes
crediticias. El agente combina tres estrategias de decisión:

\begin{enumerate}
	\item \textbf{Regla de Bayes}: minimiza el riesgo esperado condicional a los
	      datos observados.
	\item \textbf{Regla Minimax}: minimiza la peor pérdida posible sobre un
	      intervalo de credibilidad posterior.
	\item \textbf{Escalación a humano}: deriva la decisión a un analista cuando
	      la incertidumbre epistémica es demasiado alta y las reglas discrepan.
\end{enumerate}

Un \textbf{meta-selector} elige entre estas tres basándose en la entropía de
la distribución posterior y en el acuerdo entre las dos reglas automáticas.
El agente está implementado como un grafo acíclico dirigido (DAG) en
LangGraph, donde cada nodo es una función pura sobre el estado del agente.

\subsection*{Convención de pérdida (Statlog)}

La matriz de pérdida sigue la convención del dataset \emph{German Credit}
de Statlog, donde clasificar a un solicitante moroso como bueno es 5~veces
más costoso que clasificar a uno bueno como moroso:

\begin{equation}\label{eq:loss-matrix}
	L(\theta,a) =
	\begin{array}{c|ccc}
		\theta \backslash a & \text{accept} & \text{reject} & \text{escalate} \\ \hline
		\text{Pagador}         & 0             & 1             & 1.5             \\
		\text{Moroso}          & 5             & 0             & 1.5
	\end{array}
\end{equation}

La acción \texttt{escalate} tiene un costo fijo de~1.5, superior al riesgo
máximo de las reglas puras ($\max_{p\in[0,1]}\min(5p,1-p)=5/6\approx0.833$),
lo que garantiza que está \textbf{estrictamente dominada} para cualquier
$p\in[0,1]$. Solo el meta-selector puede activarla.

% ===================================================================
\section{Arquitectura del grafo (DAG)}
% ===================================================================

El agente se compone de cinco nodos funcionales encadenados secuencialmente.
Cada nodo recibe el estado completo del agente y devuelve una actualización
parcial que LangGraph fusiona automáticamente.

\begin{figure}[h]
	\centering
	\begin{tikzpicture}[
			node distance = 1.2cm and 2.8cm,
			box/.style = {
					rectangle, rounded corners, draw, very thick,
					minimum width=3.0cm, minimum height=1.2cm,
					align=center, font=\small\bfseries
				},
			edge/.style = {->, thick, >=Stealth},
			decision/.style = {
					diamond, draw, very thick, aspect=1.8,
					minimum width=3.5cm, align=center, font=\small\bfseries
				},
			terminal/.style = {
					rectangle, draw, very thick, minimum width=2.0cm,
					minimum height=0.8cm, font=\small\bfseries, rounded corners=2mm
				},
			every edge/.append style = {edge}
		]

		% Nodes
		\node[box,fill=blue!15]   (obs)      {Observar\\$x\in\X$};
		\node[box,fill=green!15]  (post)     [right=of obs] {Posterior\\$\theta\mid x \sim \Beta(\alpha,\beta)$};
		\node[box,fill=yellow!15] (bayes)    [right=of post] {Riesgo\\Bayes\\$\delta_B(x)$};
		\node[box,fill=orange!15] (minimax)  [right=of bayes] {Riesgo\\Minimax\\$\delta_M(x)$};
		\node[decision,fill=red!10] (meta)   [right=of minimax] {Meta\\Selector};
		\node[terminal]            (end)     [right=of meta] {END};

		% Edges
		\draw[edge] (obs)      -- node[above,font=\small] {$x$} (post);
		\draw[edge] (post)     -- node[above,font=\small] {$\alpha,\beta,H$} (bayes);
		\draw[edge] (bayes)    -- node[above,font=\small] {$R_B(a)$} (minimax);
		\draw[edge] (minimax)  -- node[above,font=\small] {$R_M(a)$} (meta);
		\draw[edge] (meta)     -- (end);

		% Annotations below
		\node[below=0.5cm of post,   font=\small,text width=3cm,align=center]
		{\texttt{make\_posterior\_node}};
		\node[below=0.5cm of bayes,  font=\small,text width=3cm,align=center]
		{\texttt{riesgo\_bayes}};
		\node[below=0.5cm of minimax,font=\small,text width=3cm,align=center]
		{\texttt{riesgo\_minimax}};
		\node[below=0.5cm of meta,   font=\small,text width=3cm,align=center]
		{\texttt{meta\_selector}};

		% Legend
		\node[above=0.2cm of obs, font=\small\itshape] {Nodo 1};
		\node[above=0.2cm of post, font=\small\itshape] {Nodo 2};
		\node[above=0.2cm of bayes, font=\small\itshape] {Nodo 3};
		\node[above=0.2cm of minimax, font=\small\itshape] {Nodo 4};
		\node[above=0.2cm of meta, font=\small\itshape] {Nodo 5};

	\end{tikzpicture}
	\caption{Grafo acíclico dirigido (DAG) del agente de Wald. Cada nodo es
		una función pura \texttt{AgentState $\to$ AgentState}. Las variables que
		fluyen entre nodos se indican sobre las aristas.}
	\label{fig:dag}
\end{figure}

\subsection*{Estructura del estado}

El estado del agente se define como un \texttt{TypedDict} con los siguientes
campos:

\begin{itemize}
	\item \texttt{applicant\_id}, \texttt{x}, \texttt{y\_true}: entrada.
	\item \texttt{p\_hat}, \texttt{alpha}, \texttt{beta}, \texttt{entropy\_norm}: posterior.
	\item \texttt{bayes\_risk}, \texttt{minimax\_risk}: riesgos por acción.
	\item \texttt{bayes\_action}, \texttt{minimax\_action}, \texttt{chosen\_action},
	      \texttt{rule\_used}, \texttt{realised\_loss}: decisión.
\end{itemize}

% ===================================================================
\section{Descripción de los nodos}
% ===================================================================

\subsection{Nodo 1: Observar (\texttt{observar})}

Recibe el vector de atributos $x\in\X$ del solicitante. En la implementación
actual, $x$ es inyectado directamente en el estado por el \emph{runner}
externo. El nodo valida que el campo \texttt{x} exista y retorna una
actualización vacía.

\subsection{Nodo 2: Posterior (\texttt{make\_posterior\_node})}

Convierte la salida $p = \hat{P}(\text{bad}\mid x)$ de un clasificador
(regresión logística) en una distribución Beta que representa la creencia
posterior sobre $\theta = P(\text{bad}\mid x)$:

\begin{equation}\label{eq:moment-match}
	\theta \mid x \sim \Beta(\alpha,\beta),\qquad
	\alpha = \kappa \cdot p,\quad \beta = \kappa \cdot (1-p),
\end{equation}

donde $\kappa$ es un hiperparámetro de concentración (\texttt{POSTERIOR\_KAPPA})
que controla la agudeza de la posterior. Valores grandes de $\kappa$
producen una distribución más concentrada (menos incertidumbre epistémica);
valores pequeños la ensanchan.

Este \emph{moment-matching} es una aproximación deliberada (el
\enquote{truco de Laplace--Beta}) que evita el costo de una regresión
logística Bayesiana completa, manteniendo la estructura de Wald manejable.

Además de la media $p$ y los parámetros $(\alpha,\beta)$, el nodo calcula
la entropía normalizada de la Beta:

\begin{equation}\label{eq:entropy}
	H_{\text{norm}} = \exp\bigl(H_{\Beta}(\alpha,\beta)\bigr) \in [0,1],
\end{equation}

donde $H_{\Beta}$ es la entropía diferencial de la Beta en nats.
$H_{\text{norm}} = 1$ corresponde a $\Beta(1,1)$ (máxima incertidumbre,
distribución uniforme), y $H_{\text{norm}} \to 0$ cuando la distribución
se degenera en un punto (certeza total).

\subsection{Nodo 3: Riesgo Bayesiano (\texttt{riesgo\_bayes})}

Calcula el riesgo esperado para cada acción $a\in\A$ bajo la creencia
posterior actual:

\begin{equation}\label{eq:bayes-risk}
	R_B(a) = \E_{\theta\mid x}[L(\theta,a)]
	= p \cdot L(\text{bad},a) + (1-p) \cdot L(\text{good},a)
\end{equation}

La regla de Bayes es el argumento minimizador:

\begin{equation}\label{eq:bayes-rule}
	\delta_B(x) = \operatorname*{arg\,min}_{a\in\A} R_B(a)
\end{equation}

Para la matriz de Statlog, el umbral de decisión cerrado entre aceptar y
rechazar es:

\begin{equation}\label{eq:threshold}
	\theta^* = \frac{L(\text{good},\text{reject})}
	{L(\text{good},\text{reject}) + L(\text{bad},\text{accept})}
	= \frac{1}{1+5} = \frac{1}{6} \approx 0.1667
\end{equation}

Si $p < 1/6$ la regla de Bayes acepta; si $p > 1/6$ rechaza. La acción
\texttt{escalate} nunca es seleccionada porque su costo $1.5$ excede el
riesgo mínimo de las otras dos acciones para cualquier $p$.

\subsection{Nodo 4: Riesgo Minimax (\texttt{riesgo\_minimax})}

En lugar de usar un único $p$, construye un \textbf{intervalo de credibilidad}
central al $(1-\gamma)\%$ sobre la posterior Beta y calcula el riesgo en el
peor caso:

\begin{align}
	p_{\text{lo}} & = F_{\Beta}^{-1}\bigl(\tfrac{\gamma}{2}; \alpha,\beta\bigr),   \\
	p_{\text{hi}} & = F_{\Beta}^{-1}\bigl(1-\tfrac{\gamma}{2}; \alpha,\beta\bigr), \\
	R_M(a)        & = \sup_{p\in[p_{\text{lo}},p_{\text{hi}}]}
	\bigl[p \cdot L(\text{bad},a) + (1-p) \cdot L(\text{good},a)\bigr].
\end{align}

Dado que el riesgo esperado es \textbf{lineal} en $p$, el supremo se alcanza
necesariamente en uno de los dos extremos del intervalo:

\begin{equation}\label{eq:worst-case}
	R_M(a) = \max\bigl(R_B(p_{\text{lo}};a),\; R_B(p_{\text{hi}};a)\bigr)
\end{equation}

Esto se verifica experimentalmente en el test \texttt{test\_worst\_case\_at\_endpoints}.
La regla minimax es:

\begin{equation}\label{eq:minimax-rule}
	\delta_M(x) = \operatorname*{arg\,min}_{a\in\A} R_M(a)
\end{equation}

Por construcción, $\delta_M$ es \textbf{más conservadora} que $\delta_B$:
rechaza al menos con la misma frecuencia (test~\texttt{test\_minimax\_is\_more\_conservative\_than\_bayes}).

\subsection{Nodo 5: Meta-selector (\texttt{meta\_selector})}

Selecciona entre las tres estrategias basándose en dos criterios:

\begin{enumerate}
	\item \textbf{Entropía posterior} $H_{\text{norm}}$: cuantifica la
	      incertidumbre epistémica.
	\item \textbf{Acuerdo entre reglas}: si $\delta_B(x) = \delta_M(x)$ o no.
\end{enumerate}

La lógica de decisión es la siguiente:

\begin{figure}[h]
	\centering
	\begin{tikzpicture}[
			node distance = 1.0cm and 2.5cm,
			dec/.style = {diamond, draw, thick, aspect=1.8,
					minimum width=2.8cm, align=center, font=\small},
			act/.style = {rectangle, draw, thick, rounded corners,
					minimum width=2.2cm, minimum height=0.9cm,
					align=center, font=\small\bfseries},
			edge/.style = {->, thick, >=Stealth}
		]

		\node (start)   [act,fill=gray!10]                        {Estado\\$H, \delta_B, \delta_M$};
		\node (dec1)    [dec,fill=red!10,below=of start]          {$H \ge \tau_{\text{esc}}$\\$\land\ \delta_B \neq \delta_M$?};
		\node (dec2)    [dec,fill=orange!10,below=of dec1]        {$H \ge \tau_{\text{mm}}$?};
		\node (bayes)   [act,fill=green!15,below left=of dec2]    {Usar\\$\delta_B(x)$};
		\node (minimax) [act,fill=blue!15,below right=of dec2]    {Usar\\$\delta_M(x)$};
		\node (esc)     [act,fill=red!20,right=of dec1]           {Escalar\\a humano};

		\draw[edge] (start)  -- (dec1);
		\draw[edge] (dec1)   -- node[right,font=\small] {Sí} (esc);
		\draw[edge] (dec1)   -- node[right,font=\small] {No} (dec2);
		\draw[edge] (dec2)   -- node[above,font=\small] {No} (bayes);
		\draw[edge] (dec2)   -- node[above,font=\small] {Sí} (minimax);

	\end{tikzpicture}
	\caption{Árbol de decisión del meta-selector. $\tau_{\text{esc}} = 0.65$,
		$\tau_{\text{mm}} = 0.45$ con los parámetros actuales.}
	\label{fig:meta}
\end{figure}

\begin{itemize}
	\item \textbf{Entropía baja} ($H < \tau_{\text{mm}}$): se confía en la
	      regla de Bayes (eficiente).
	\item \textbf{Entropía alta} ($H \ge \tau_{\text{mm}}$): se usa la regla
	      minimax (robusta).
	\item \textbf{Entropía muy alta y desacuerdo} ($H \ge \tau_{\text{esc}}$
	      y $\delta_B \neq \delta_M$): se escala a un analista humano.
\end{itemize}

Los umbrales $\tau_{\text{mm}}$ y $\tau_{\text{esc}}$ están definidos en
\texttt{config.py} y son editables.

% ===================================================================
\section{Verificación empírica}
% ===================================================================

\subsection*{Configuración experimental}

\begin{itemize}
	\item Dataset: German Credit (Statlog), 1000 filas, 20 atributos, 31.4\% morosos.
	\item Clasificador base: regresión logística con pipeline
	      (StandardScaler + OneHotEncoder).
	\item Split: 80/20 estratificado, semilla~42.
	\item Concentración posterior: $\kappa = 5$ (baja, para activar el meta-selector).
\end{itemize}

\subsection*{Resultados}

\begin{table}[h]
	\centering
	\begin{tabular}{lrc}
		\toprule
		\textbf{Métrica}      & \textbf{Valor}       \\ \midrule
		Precisión (train)     & 0.748                \\
		Precisión (test)      & 0.740                \\
		Umbral Bayes cerrado  & $1/6 \approx 0.1667$ \\
		$p_{\max}$ aceptados  & 0.1665               \\
		$p_{\min}$ rechazados & 0.1667               \\ \bottomrule
	\end{tabular}
	\caption{La coincidencia entre el umbral cerrado y el observado
		($0.1665/0.1667$) verifica empíricamente el Teorema de Bayes.}
	\label{tab:threshold}
\end{table}

\begin{table}[h]
	\centering
	\begin{tabular}{lrc}
		\toprule
		\textbf{Regla} & \textbf{Pérdida total} \\ \midrule
		Always-accept  & 315.00                 \\
		Always-reject  & 137.00                 \\
		Bayes puro     & 140.00                 \\
		Minimax puro   & 134.00                 \\
		Meta-selector  & 140.00                 \\ \bottomrule
	\end{tabular}
	\caption{Pérdida total realizada sobre 200 solicitudes de test.
		Minimax supera a Bayes en esta ejecución (efecto de muestra finita).}
	\label{tab:loss}
\end{table}

\begin{table}[h]
	\centering
	\begin{tabular}{lcc}
		\toprule
		               & \multicolumn{2}{c}{\textbf{Minimax}}          \\
		\textbf{Bayes} & accept                               & reject \\ \midrule
		accept         & 9                                    & 60     \\
		reject         & 0                                    & 131    \\ \bottomrule
	\end{tabular}
	\caption{Tabla de contingencia entre reglas. En 60 casos (30\%)
		discrepan: Bayes acepta donde minimax rechaza.}
	\label{tab:confusion}
\end{table}

\subsection*{Distribución de entropía}

Con $\kappa=5$, la entropía normalizada tiene media $0.52$ y alcanza hasta
$0.82$. El meta-selector asigna 129 casos a minimax y 71 a Bayes. La ruta
de escalación no se activa porque las condiciones de alta entropía y
desacuerdo entre reglas están \textbf{anti-correlacionadas}: la entropía
es máxima cuando $p\approx 0.5$ (ambas reglas coinciden en rechazar), y
el desacuerdo ocurre cerca del umbral $p\approx 1/6$ donde la entropía es
moderada.

% ===================================================================
\section{Validación teórica (tests)}
% ===================================================================

Los 7 tests de cordura verifican propiedades fundamentales:

\begin{enumerate}
	\item \textbf{Umbral de Bayes}: coincide con $1/6$ (test 1).
	\item \textbf{Frontera de decisión}: $p < 1/6 \to$ accept, $p > 1/6 \to$ reject (test 2).
	\item \textbf{Minimax es más conservador}: rechaza al menos tantos como Bayes (test 3).
	\item \textbf{Escalación bajo desacuerdo + alta entropía}: el meta-selector escala (test 4).
	\item \textbf{Bayes bajo certeza}: con entropía baja, usa Bayes (test 5).
	\item \textbf{Linealidad del riesgo}: $\E[L\mid p,a]$ es lineal en $p$ (test 6).
	\item \textbf{Peor caso en extremos}: $\sup_{[p_{\text{lo}},p_{\text{hi}}]}$
	      se alcanza en un extremo (test 7).
\end{enumerate}

% ===================================================================
\section{Configuración y parametrización}
% ===================================================================

Todos los parámetros editables residen en \texttt{config.py}:

\begin{table}[h]
	\centering
	\begin{tabular}{lll}
		\toprule
		\textbf{Parámetro}                      & \textbf{Default} & \textbf{Efecto}          \\ \midrule
		\texttt{SEED}                           & 42               & Reproducibilidad         \\
		\texttt{POSTERIOR\_KAPPA}               & 5.0              & Concentración Beta       \\
		\texttt{ENTROPY\_THRESHOLD}             & 0.45             & Activación de minimax    \\
		\texttt{ESCALATE\_ENTROPY\_THRESHOLD}   & 0.65             & Activación de escalación \\
		\texttt{MINIMAX\_CRED\_LEVEL}           & 0.90             & Amplitud del intervalo   \\
		\texttt{LOSS{[}(``bad'',``accept''){]}} & 5.0              & Asimetría 5:1            \\
		\texttt{TEST\_SIZE}                     & 0.20             & Proporción de test       \\ \bottomrule
	\end{tabular}
	\caption{Parámetros editables del agente.}
	\label{tab:params}
\end{table}

\section*{Conexión con el Capítulo 8 (MDPs)}

Cada decisión individual constituye un MDP de un paso donde
$\Risk(\theta,\delta) = -V^\pi(s)$. El isomorfismo Wald--MDP del Capítulo~11
se manifiesta al extender este agente a decisiones secuenciales
(préstamos renovables con observación posterior del desempeño del prestatario).

\end{document}
